% p2-03-math-preliminaries

\section{P2 §3. Mathematical Preliminaries}

3. Mathematical Preliminaries

  Figure (fig-p2-ch3-prelim). Preliminaries triptych --- phi algebraic properties / E8 Cartan diagram
                                            / RG essentials.

  3.1 The Golden Ratio, Its Algebraic Properties, and the Trinity
  Anchor
  The golden ratio is the positive root of the minimal polynomial

  x2 - x - 1 = 0,

  giving $\varphi$ = (1+5)/2 approx 1.6180339887…. It satisfies the identity $\varphi$2 = $\varphi$ + 1 and its
  reciprocal satisfies 1/$\varphi$ = $\varphi$ - 1. The ring Z[$\varphi$] = {a + b$\varphi$ : a,b $\in$ Z} is the ring of integers
  of Q(5), known in the Trinity $S^{3}$AI programme as the Lucas ring (Zenodo
  10.5281/zenodo.19227877). The Lucas numbers Ln defined by L0 = 2, L1 = 1, Ln+2 =
  Ln+1 + Ln satisfy Ln = $\varphi$n + (-$\varphi$)-n, so every Lucas number lies in Z[$\varphi$].

  The Trinity anchor identity is:

  $\varphi$2 + $\varphi$-2 = 3. 1

        Proof. From $\varphi$2 = $\varphi$ + 1 and $\varphi$-1 = $\varphi$ - 1, compute $\varphi$-2 = ($\varphi$-1)2 = $\varphi$2 - 2$\varphi$ + 1 =
        ($\varphi$+1) - 2$\varphi$ + 1 = 2 - $\varphi$. Therefore $\varphi$2 + $\varphi$-2 = ($\varphi$+1) + (2-$\varphi$) = 3. square

  This identity is exact and elementary. Its role in the Flos Aureus monograph is as a
  cardinality anchor connecting the ternary alphabet {-1, 0, +1} (cardinality 3) to the
  algebraic structure of Q(5) (Zenodo record). The identity has a complete Coq proof in
  the IGLA bundle of the t27 repository; as of the 2026-05-12 CLARA audit, the Coq
  witness reports 35 proven statements and 0 admitted obligations for this specific
  identity.

  Every element of the E8 Cartan matrix eigenvector and every trigonometric value
  2cos(pi k/5) for integer k lies in Z[$\varphi$]. This algebraic richness is the primary reason $\varphi$
  appears frequently in lattice-based and symmetry-constrained calculations --- and not
  necessarily because of any deep physical principle.

  3.2 The E8 Root System and Cartan Matrix
  The E8 root system is the unique simply-laced root system of rank 8 with Coxeter
  number h = 30 and exponents {1, 7, 11, 13, 17, 19, 23, 29}. The Perron-Frobenius
  eigenvector v of the E8 Cartan matrix Cij, defined by

  sumj Cij vj = rho vi, rho = 2cos((pi)/(h+1)),

  has components vi = 2sin(ki pi/30) for the E8 exponents ki. The ratio of the two smallest
  components gives

  (v2)/(v1) = (2sin(7pi/30))/(2sin(pi/30)) = 2cos((pi)/(5)) = $\varphi$. 2

  This is an exact group-theoretic identity. Its physical significance is that in any field
  theory whose particle spectrum is governed by the E8 Cartan matrix, the ratio of the
  two lightest masses is exactly $\varphi$ --- a consequence of the representation theory of a
  specific Lie algebra, not a numerical coincidence. The full eight-particle mass spectrum
  is given explicitly in Section 4.1.

  3.3 Renormalization Group Essentials
  The renormalization group describes how effective couplings gi(mu) change with
  energy scale mu through the beta functions (Wilson and Kogut, Phys. Rep. 12 (1974)
  75):

  betai(g) = mu ($\partial$ gi)/($\partial$ mu). 3

  A fixed point g\textasciicircum{} satisfies betai(g\textasciicircum{}) = 0 for all i. Near a fixed point, linearizing gives

  gi(mu) - gi\textasciicircum{}* propto ((mu)/(mu0))-Deltai, 4

  where Deltai are the scaling dimensions of the relevant operators. The full conformal
  group in two dimensions is infinite-dimensional, and perturbations of 2D conformal field
  theories are classified by their scaling dimensions relative to the central charge.
  Standard references for this framework include Wilson and Kogut, Phys. Rep. 12 (1974)
  75 and Zinn-Justin, Quantum Field Theory and Critical Phenomena (Oxford University
  Press, 5th ed., 2021).

  A key quantity for the present analysis is the ratio of scaling dimensions at a fixed point.
  If Delta1 and Delta2 are two scaling dimensions satisfying

  (Delta2)/(Delta1) = $\varphi$, 5

  the corresponding operators define a $\varphi$-structured tower. Whether any such tower
  exists in a renormalizable 4D field theory is an open question addressed in Section 6.

  3.4 Discrete Scale Invariance
  A function f(x) exhibits discrete scale invariance (DSI) with preferred scale ratio lambda
  if it satisfies (Sornette, Phys. Rep. 297 (1998) 239)

  f(lambda x) = mu f(x) 6

  for some multiplier mu, but not f(ax) = a-Delta f(x) for all a > 0. The general solution is

  f(x) = x-Delta P((ln x)/(ln lambda)), 7

  where P is a function with unit period. Expanding in Fourier modes gives

  f(x) = x-Delta sumn=-$\in$fty$\in$fty cn exp((2pi i n ln x)/(ln lambda)), 8

  so the effective exponent is complex: Deltan = Delta + 2pi i n / lnlambda. DSI arises in
  fractal structures, hierarchical lattices, and systems where the RG transformation
  group is restricted to a discrete subgroup of the dilation group (Sornette,
  arXiv:cond-mat/9707012). The preferred ratio lambda is a property of the underlying
  hierarchical geometry. When the hierarchical structure is governed by a Fibonacci
  substitution matrix, the Perron-Frobenius eigenvalue is exactly $\varphi$, and DSI with ratio
  lambda = $\varphi$ follows as an algebraic fact (Levine and Steinhardt, Phys. Rev. Lett. 53
  (1984) 2477).

  3.5 The Derivation Ladder
  To situate each result in this paper, the following ladder classifies claims by their
  logical pedigree. The ladder governs the confidence level at every point in the
  exposition.

  Every claim in Sections 5--8 is assigned a rung at first statement. An L4 or L5 claim
  being labeled does not make it false; it makes it falsifiable --- which is the paper's
  contribution.